\documentclass{report}
\usepackage{amsmath,amssymb}
\setlength{\parindent}{0mm}
\setlength{\parskip}{1em}
\begin{document}
\begin{center}
\rule{6in}{1pt} \
{\large Alex Gorodetsky \\
{\bf Function-train: a continuous analogue to the tensor-train decomposition}}

Massachusetts Institute of Technology \\ 77 Massachusetts Avenue \\ Room 37-312 \\ Cambridge \\ MA 02139
\\
{\tt goroda@mit.edu}\\
Sertac Karaman\\
Youssef Marzouk\end{center}

The curse of dimensionality often arises in high-dimensional uncertainty
quantification problems. For example, solving a stochastic partial
differential equation (PDE) with tens to hundreds of independent
stochastic inputs can be prohibitively expensive. Recent work on the
factorization of high-dimensional arrays has shown great potential for
alleviating this computational burden when problems exhibit low-rank
structure. Such factorizations, termed tensor decompositions, often allow
for the solution of PDEs in time that scales linearly with dimension and
polynomially with rank. Furthermore, they enable fast post-processing of
various quantities of interest through multilinear algebra methods that
also scale linearly with dimension and polynomially with rank.

In this talk we extend a particular tensor decomposition, the
tensor-train, to the continuous case of multidimensional functions. This
extension allows us to easily incorporate discretization adaptivity in
our approximation by breaking away from the typical grids associated with
tensor decompositions. Furthermore, the continuous tensor decomposition
-- which we term the "function-train"-- provides a framework for
post-processing UQ-related quantities that are discretized at different
levels, thereby increasing the flexibility and applicability of these
algorithms.

We begin by describing a new cross-approximation algorithm for computing
the CUR/skeleton decomposition of bivariate functions. We then extend
this decomposition to the multidimensional case of the function-train
decomposition. Computing the decomposition relies on continuous analogues
of matrix factorizations, such as continuous QR and LU factorizations of
matrix-valued functions. We continue by describing a continuous, or
functional, alternating least squares optimization algorithm for fast
multilinear algebraic operations such as multiplication of low-rank
functions and application of low-rank operators. The computational cost
of these operations scales linearly with dimension and polynomially with
rank. Finally, we apply our methods to UQ-related problems such as
Gaussian filtering and the solution of stochastic elliptic partial
differential equations.


\end{document}
